\documentclass{article}

\usepackage[T1]{fontenc} \usepackage{amsmath} \usepackage{amsthm}
\usepackage{amssymb} \usepackage{graphicx} \usepackage{float}

%%% OUTLINES
\newtheorem{theorem}{Theorem}[section] \newtheorem{lemma}[theorem]{Lemma}
\newtheorem{corollary}[theorem]{Corollary}

\theoremstyle{definition} \newtheorem{definition}[theorem]{Definition}
\newtheorem{example}[theorem]{Example}

\theoremstyle{remark} \newtheorem{remark}[theorem]{Remark}
%%% END OUTLINES

%%% SYMBOLS
\newcommand{\R}{\mathbb{R}}  % The real numbers. 
\newcommand{\E }{\mathbb{E}}  %Expectation Operator \renewcommand{\P}{\mathbb{P}}  % Probability
\newcommand{\Cov}{\operatorname{Cov}}  % Cov Operator
\newcommand{\Var}{\operatorname{Var}}  % Var Operator
\newcommand{\MSEPin}{\textbf{MSEP}_{\text{in}}} \newcommand{\MSE}{\textbf{MSE}}
\newcommand{\eqdis}{\stackrel{d}{=}}

\newcommand{\bigzero}{\mbox{\normalfont\Large\bfseries 0}}
\newcommand{\bigone}{\mbox{\normalfont\Large\bfseries I}}



%%% END SYMBOLS

\title{Notes on Trees} \author{Tree Team}

\begin{document} \maketitle

%\begin{abstract} Abstract \end{abstract}

\section{Introduction} 

A popular model for supervised learning is the Classification And Regression
Tree (CART). As in all of machine learning, question regarding how well a model
generalizes to unseen data is central. The problem of model generalzation is
sometimes called the overfitting problem. One way of evaluating generlization is
using cross-validation, which allocates a prespecified fraction of the data for
evaluating the trained model. This type of validation has many caveats, for one,
there is less data available for training the model. Further, using too little
data for validation produces poor estimate of the \textbf{generalzation error}.
Others have proposed estimating the error using the same data that generated the
model using various methods (AIC, BIC etc). 

Finally, another approach is to evaluate the generalization by using inference
based methods and hypothesis testing. The connection between direct estimation
of generalization error and hypothesis test is unclear. However, they aim to do
the same thing. 

For CART there are a multitude of methods propsed for dealing with overftting
using inference. \cite{shih2004variable} use inference to determine the split of
a node. Others use traditional scoring methods for choosing splits but use
inference to determine the quality of the split, see
\cite{jensen1997adjusting, hothorn2006unbiased, neufeld2022tree}. However, care
should be taken when using inference. \cite{jensen1997adjusting} argues that
mutiple testing across covariates should be accounted for and
\cite{neufeld2022tree} warns about \textit{selective inference}. See
\cite{neufeld2022tree} for selective inference on CART and
\cite{fithian2017optimal} for the general theory of selective inference. 

Results from the change-point detection literature can be leveraged for
constructin trees, since splitting a node is essentially the same as finding a
change-point in the data. Specifically, we will use results from
\cite{hawkins1977testing, yao1986asymptotic} for inference on trees. There are a
multitude of ways for constructing trees using inference, but, below we will
focus on a few methods that we argue are reasonable. Furthermore, we will
analyse the effect of inference based methods for boosted trees. 

\section{Background}

\subsection{Inference}

There are in general two broad methods for constructing CART using
inference; constructing trees recursively using inference (top-down ) and pruning trees
using inference (bottomn-up). One can also combine both methods to construct trees.

Constructing trees recursively, can be done by evaluating the quality of a split
by some suitable hypothesis test. For instance, let $ \mathcal{D} = \left\{
(X_i,y_i) \right\} _{i=1}^{n} $ be data available at a particular node of
interest. The question is: can $ \mathcal{D} $ be partioned into two groups $
\mathcal{D}_A, \mathcal{D}_B $ in a statistically significant way. Typically,
these groups are chosen by minimizing some loss $ l(\mathcal{D}; A, B) =
R(\mathcal{D}_A) + R(\mathcal{D}_B) $, for some function $ R $, e.g the $ L^2 $
loss $ R(\mathcal{D}) = \sum_{i=1}^{n}(y_i - \bar{y})^2$. Then, a hypothesis
test on the estimated partion is performed. The null hypothesis being that there
is no difference between the groups. E.g $ H_0: \mu_A = \mu_B $
\cite{neufeld2022tree}, both groups have the same mean. Then a suitable test is
performed, and the split is accepted if $ H_0 $ is rejected with a significance
$ \alpha $. This continues recursively, until no more splits can be made.

On the other hand, constructing trees by pruning is done by first building a
full tree, for instance by minimizing the loss $ l $ described above in a greedy
manner. Once a tree is fully constructed, one can for instance, evaluate the quality
of splits in the \textit{frotier nodes}(nodes with children that are leaves) by
performing a similar hypothesis test as above and removing splits that are
insignificant. Finally, the pruning can be stopped once no more leaves can be
pruned, this is done by \cite{jensen1997adjusting}.  

In both methods of constructing trees, we are in some sense performing a test
for each split, this puts us in the multiple testing regime. Here, we can not
view each test as being performed independently, in some cases, a test is
performed only if a previous test is rejected. However, understanding how these
events relate to each other may be too complicated. Instead, we can bound the
significance of all test together using \textit{Bonfferoni} corrections. More
precisely, we want to bound the \textit{type-1} error in the following way:

\begin{align*}
  &\mathbb{P}(\textrm{rejecting some } H_0^{(i)} | H_0^{(i)} \textrm{ is true for
  some }i) = \\
  = &\mathbb{P}(\cup_i \{  \textrm{Reject } H_0^{(i)} \} | H_0^{(i)} \textrm{ is
  true for some $ i $}) = \\ 
  \leq  & \sum_{i=1}^{N}\mathbb {P}(\{ p_i \leq \alpha \} | H_0^{(i)} 
  \textrm{ is true} ) \leq N\alpha \leq \tilde{\alpha}.
\end{align*}

Where, $ N $ is the number of nodes in the tree and $ \tilde{\alpha} $ measures
the total uncertainty. A natural choice for $ \alpha $  is then $ \alpha =
\tilde{\alpha}/N $. However, if the number of nodes $ N $ in the tree is not
deterministic, this choice of $ \alpha $ is not suitable. 

Instead, summing the p-values $ p_i $ for each node of the tree can be used as $
\alpha $. In the spririt of the top-down and bottomn-up, this means that the
procedure stops just before or after the condition $ \sum_{i=1}^{} p_i \leq
\tilde{\alpha } $ is satisfied. This is a reasonable stopping rule for
constructing trees. A potential issue with this approach is the order in which
we consider splits. For instance, building a tree top-down, at each recursion
step, there will be $ M $ splits to consider, in which order should the p-values
for these splits be added to the total? Further, say we accept a split, should
the subsequent split be considerd as well? There is no best way to do this. A
similar problem arises for pruning. Therefore, a \textit{splitting sequence}
needs to be specified. 

If we can consider \textit{cost complexty pruning} we can reconcile this
problem, since the algorithm specifies a splitting sequence. 

\cite{jensen1997adjusting} proposes the idea of considering multiple testing
within a split when there are multiple covariates considered. In fact, a split
is determined by comparison to other covariates. For each covariate, a best
split is determined, then the chosen split is the best among these individual
choices. In this case, a multiple test is being performed. If we consider this,
then we can view our problem of constructing a tree as a multiple-multiple test.
Should we once again use a Bonfferoni bound $ \alpha = \tilde{\alpha}/pN $,
where $ p $ is the number of covariates and $ N $ is the number of nodes. If
there are many covariates, will this be too restrictive?

There are many configurations we could consider, however, we will focus on the
following:
\begin{itemize}
  \item Expand this more
  \item Recursive splitting + summing the smallest p-value.
  \item ccp-pruning  + summing p-values.
  \item boosting with fixed learner size + summing p-values. 
  \item Adaptive(building with recursion and sum limit) boosting + summing p-values for trees. 
\end{itemize}

Now that we have determined the configurations that should be considered, we can
focus on how we should construct the hypothesis test. We have tested creating a
test statisic using an estimate of $ MSEP $ however we see that this results in
some sort of selective inference and since we want to avoid selective inferece
all togehter, we have chosen to go ahead with the test proposed by
\cite{hawkins1977testing}. We have not decided yet wheter this test-statistic is
free of selective inference but that is the conjecture. 

\subsection{CART}

Below, are some descriptions on various algorithms that may be used in this
project. 

\subsubsection{Cost-Complexity Pruning}

\label{sec:ccp}

Breiman's method of minimal cost-complexity builds a full tree and sequencially
prunes nodes to construct a tree that is less complex and that is supposed to
generalize better. 

If we denote $ T $ as a tree and define $ R(T) $ as the weighted loss
contributions from the leaves of $ T $, we can define a penalized loss $
R_{\alpha}(T) = R(T) + \alpha|T|$, where $ |T| $ is the number of leaves of the
tree $ T $ and $ \alpha $ is some hyper-parameter that penalizes complexity.
Breiman's algorithm produces the best \textbf{subtree} for a given $ \alpha $
that minimizes $ R_{\alpha}(T) $. Furthermore, it can be shown that the best
trees are nested subtrees $ T_k $ of $ T $ such that each is optimal for an
interval of $ \alpha $. That is, there exits a sequence $ \{\alpha_k \} $ 

\begin{align}
 0 = \alpha_0 < \alpha_1 < \cdots < \alpha_m < \infty 
\end{align}

such that $ T_i $ is the optimal  tree w.r.t $ R_{\alpha}(\cdot) $ for $ \alpha \in [\alpha_i,
\alpha_{i+1}) $. Here it should be clear that $ T_0 = T $ and $ T_m $ is the
trivial tree. What this essentially means is that, given some value $ \alpha \in [\alpha_i,
\alpha_{i+1}) $ the tree that minmizes $ R_{\alpha}(\cdot) $ is $ T_i $ and $ T_{i+1} $ is a subtree of $
T_i $, so starting at $ \alpha = 0 $ and increasing $ \alpha $ will produce a
sequence of trees that are pruned subtrees of the preceding until the trivial tree
$ T_m $ is reached.

The procedure for attaining $ \{\alpha_k \} $ and $ T_k $ is as follows. Given a non-trivial
tree $ T_i $, denote all non-terminal nodes with some $ t $ and calculate the 
gain $ g(t, T^{(t)}) = \frac{R(t) - R(T^{(t)})}{|T^{(t)}| - 1} $ for all $ t $, where $ R(t) $ is
the loss at the node $ t $ treated as a trivial tree and $ R(T^{(t)}) $ is the
loss of the tree that is rooted at $ t $. Then $\alpha_i = \underset{t}\min \; 
g(t,T^{(t)})$. Pruning all nodes where the $ g(\cdot) \geq  \alpha_i $ produces a new
tree $ T_{i+1} $. Now perform the same procedure on $ T_{i+1} $ and so on until
the trivial tree is reached. This results in a sequence $ \{\alpha_k \} $ and a
sequnce of nested trees $ \{ T_k \} $. Two remarks can be made. Firstly, the procedure
handles ties in gain by choosing the smallest tree as next in line. Second, the
gain is attained from simple algebra by solving for $ \alpha $ in the equation
$R_{\alpha}(t) - R_{\alpha}(T^{(t)}) = 0 $.


\subsection{Boosting}

\subsubsection{Adaptive Boosting}

When boosting regression trees in the traditional sense, there are two
parameters that need to be decided. First, the complexity of the weak learners,
e.g. number of leaves of the trees, at each boosting step and second the
stopping criterion, that is, step $ k $ to stop boosting. The complexity of the
tree is choosen heuristically. However, the stopping criterion is typically
decided using cross validation. Where the training is stopped as soon as the
out-of-sample error starts increasing. Adaptive Boosting Trees (ABT) tackles this
problem by proposing an algorithm for boosting with a built-in stopping rule
that is based on the complexity of the weak learner. That is, it is sufficient
to choose a parameter $ J $ and the procedure will produce boosted trees with at
most complexity $ J $ and it is conjectured to stop in finite time. 

ABT leverages Breiman's pruning method (see \ref{sec:ccp}) to produce learners
of at most $ J $ leaves. That is, at each boosting iteration a tree is fitted by
choosing the subtree $ T_k $ of $ T $ such that 

\begin{align*}
  \label{eq:abtrule} 
  |T_{k+1}| > J \textrm{ and } |T_{k}| \leq J.
\end{align*}

If $ |T_k| = 1$ then the procedure stops. The reason for stopping when this
criterion is met, is that the intercept will be identically zero, thus, boosting
after this criterion is met will further produce intercepts of zero which is
unecessary. Note that in the ABT paper they present a general framework for
choosing the response that should be fit at each iteration given a deviance
loss. However when an $ L^2 $ loss is considered, the reponse at each iteration
is the residuals.

A problem raised in the ABT paper is that their procedure sometimes stops early,
"stuck in local optima". They propose using a bagging fraction $ \beta $ to
overcome this issue.

\subsubsection{P-Boosting}
Adaptive or non adaptive boosting using inference.

\bibliographystyle{apalike}
\bibliography{refs}

\end{document}
